\section{Introduction}
\label{sec:intro}

Vlogs represent a unique form of blogging, distinguished by their utilization of video as the primary medium rather than text.
\blfootnote{$\spadesuit$ Interns at Shanghai AI Laboratory. $\heartsuit$ Corresponding authors.}
% Vlog is a form of blog where the medium is video instead of words.
Due to more lively expression in the dynamic scenes,
vlog has become one of the most popular online-sharing ways in the digital world. 
In the past two years,
the remarkable success of diffusion models \cite{song2021scorebased, ho2020denoising, ho2021classifierfree} has shown a great impact on video creation \cite{singer2023makeavideo, Ho2022ImagenVH, blattmann2023align, Ge2023PreserveYO, Wang2023ModelScopeTT, Yin2023NUWAXLDO, Zhang2023Show1MP, Wang2023LAVIEHV}.
Hence,
there is a natural question, 
\textit{can we build a generic AI system to generate wonderful vlogs automatically?}


Regrettably, the majority of current video diffusion approaches mainly generate short videos with a few seconds, by temporal adaptation of image diffusion models \cite{singer2023makeavideo, blattmann2023align, Wang2023ModelScopeTT, Zhang2023Show1MP, Wang2023LAVIEHV, Ge2023PreserveYO}.
In contrast,
a vlog typically constitutes a \textit{minute}-level long video in the open world. %with diversified content
Recently,
there have been some attempts in long video generation \cite{villegas2023phenaki, Yin2023NUWAXLDO}.
However,
these early works either require extensive training on large well-captioned long video datasets \cite{Yin2023NUWAXLDO}
or suffer from noticeable incoherence of shot changes \cite{villegas2023phenaki}. 
Hence,
it is still challenging to generate a minute-level vlog with complex narratives and multiple scenes.


Alternatively,
we notice that
a successful vlog production is a systematical work in our realistic world,
where
the key staffs are involved in
script creation,
actor design, 
video shooting, and 
editing \cite{gao2010vlogging}. 
Drawing inspiration from this,
we believe that,
generating a long-form video blog requires elaborate systematical planning and shooting processes,
rather than simply designing a generative model.

Hence,
we propose a generic AI system for vlog generation in this paper,
namely \textbf{Vlogger},
which can smartly address this difficult task by mimicking vlog professionals with various foundation models in the core steps.
As shown in Fig. \ref{fig:teaser},
we first hire a Large Language Model (LLM) as \textit{Director} (e.g., GPT-4 \cite{OpenAI2023GPT4TR}),
due to its great power of linguistic knowledge understanding.
Given a user story,
this director schedules a four-step plan for vlog generation.
(1) \textit{Script}.
First,
we introduce a progressive script creation paradigm for the LLM Director.
By the coarse-to-fine instructions in this paradigm,
the LLM Director can effectively convert a user story into a script,
which sufficiently describes the story by a number of shooting scenes and their corresponding shooting duration.
(2) \textit{Actor}.
After creating the script,
the LLM Director reads it again to summarize the actors,
and then invokes a character designer (e.g., SD-XL \cite{Podell2023SDXLIL}) to generate the reference images of these actors in the vlog.
(3) \textit{ShowMaker}.
With the guidance of script texts and actor images,
we develop a novel ShowMaker as a videographer,
which can effectively generate a controllable-duration snippet for each shooting scene, with spatial-temporal coherence. 
(4) \textit{Voicer}.
Finally,
the LLM Director invokes a voicer (e.g., Bark \cite{sunoai2023bark}) to dub the vlog with script subtitles.


It should be noted that
our Vlogger overcomes the challenges previously encountered in long video generation tasks.
On one hand, 
it elegantly decomposes the user story into a number of shooting scenes
and designs actor images that can participate in different scenes in the vlog.
In this case,
it can reduce spatial-temporal incoherence of abrupt shot changes,
with explicit guidance of scene texts and actor images.
On the other hand, 
Vlogger crafts individual video snippets for every scene and seamlessly integrates them into a single cohesive vlog. 
Consequently, 
this bypasses the tedious training process with large-scale long video datasets. 
Via such collaboration between top-down planning and bottom-up shooting, 
Vlogger can effectively transform an open-world story into a minute-long vlog.


Furthermore, 
we would like to emphasize that,
\textbf{ShowMaker} is a distinctive video diffusion model designed for generating video snippets of each shooting scene.
From the structural perspective,
we introduce a novel Spatial-Temporal Enhanced Block (STEB) in this model.
This block can adaptively leverage scene descriptions and actor images as textual and visual prompts,
which attentively guide ShowMaker to enhance spatial-temporal coherence of script and actors.
From the training perspective,
we develop a probabilistic mode selection mechanism,
which can boost the capacity of ShowMaker by mixed training of Text-to-Video (T2V) generation and prediction.
More notably,
by sequential combination of generation and prediction mode in the inference stage,
ShowMaker can produce a video snippet with a controllable duration.
This allows our Vlogger to generate a vlog with a preferable duration,
according to the planning of each scene in the script by LLM director.

Finally,
our method achieves the state-of-the-art performance on both zero-shot T2V generation and prediction,
by expensive experiments within the popular video benchmarks.
More importantly,
our Vlogger outperforms the well-known long video generation method,
\textit{i.e.}, 
Phenaki \cite{villegas2023phenaki},
despite utilizing only 66.7\% of the training videos.
Remarkably,
our Vlogger is capable of generating over 5-minute vlogs, 
without loss of script and actor coherence in the video.




\begin{figure*}[t]
    \centering
    % \vspace{-0.3cm}
    \includegraphics[width=1\linewidth]{figs/director.pdf}
    \vspace{-0.7cm}
    \caption{
    \textbf{Top-Down Planning.} 
    Through four rounds of dialogue with the LLM, we gradually convert the user story into the final script. 
    Based on this script, we further extract actor reference images,
    and then determine which actor would star in each script scene.
    }
    \label{fig:director}
    \vspace{-0.3cm}
\end{figure*}